\paragraph{\texorpdfstring{$(y,\delta)$}{(y,delta)} 平面闭包模型：五节点、伴随二次与谱坐标反演}
沿用结论 \ref{con:xi-terminal-zm-leyang-abel-jacobi-ode} 的椭圆曲线
\[
E:\quad Y^{2}=X^{3}-X+\frac14,
\qquad
y:=X^{2}+Y-\frac12\in\QQ(E),
\qquad
\omega:=\frac{\dd X}{2Y}.
\]
令 Abel--Jacobi 坐标 \(u\) 满足 \(\dd u=\omega\)，并令
\[
\delta:=\frac{\dd y}{\dd u}\in\QQ(E).
\]
由 \(\dd y=(4XY+3X^{2}-1)\,\omega\)（结论 \ref{con:xi-terminal-zm-leyang-abel-jacobi-ode} 的证明中已给出）可得
\[
\delta=4XY+3X^{2}-1.
\]

\begin{theorem}[二阶射流的三次代数闭式]\label{thm:xi-terminal-zm-leyang-second-jet-cubic-closure}
\leanverified{paper\_xi\_terminal\_zm\_leyang\_second\_jet\_cubic\_closure}
令 \(u\) 为满足 \(\dd u=\omega\) 的 Abel--Jacobi 坐标，并令
\[
\delta=\frac{\dd y}{\dd u},
\qquad
\delta_{2}:=\frac{\dd^{2}y}{\dd u^{2}}=\frac{\dd\delta}{\dd u}.
\]
则在 \(E\) 上有显式表达式
\[
\delta_{2}=20X^{3}-12X+2+12XY.
\]
并且三元组 \((y,\delta,\delta_2)\) 满足一条三次代数关系：
\[
\begin{aligned}
0={}&64\delta^{3}+96\delta^{2}\delta_{2}+576\delta^{2}y-288\delta^{2}+48\delta\delta_{2}^{2}-144\delta\delta_{2}y-216\delta\delta_{2}\\
&+9600\delta y^{3}-5424\delta y^{2}+1200\delta y+132\delta+8\delta_{2}^{3}-216\delta_{2}^{2}y-63\delta_{2}^{2}\\
&-3200\delta_{2}y^{3}-1392\delta_{2}y^{2}+1008\delta_{2}y+128\delta_{2}+16000y^{3}-2640y^{2}-816y-16.
\end{aligned}
\]
\end{theorem}
\begin{proof}
由 \(E:Y^2=X^3-X+\tfrac14\) 得 \(\frac{\dd X}{\dd u}=2Y\) 且 \(\frac{\dd Y}{\dd u}=3X^2-1\)。
对 \(\delta=4XY+3X^2-1\) 求导得
\[
\delta_2
=4\left(\frac{\dd X}{\dd u}Y+X\frac{\dd Y}{\dd u}\right)+6X\frac{\dd X}{\dd u}
=8Y^2+12X^3-4X+12XY.
\]
再用 \(8Y^2=8(X^3-X+\tfrac14)=8X^3-8X+2\) 即得 \(\delta_2=20X^3-12X+2+12XY\)。
\par
另一方面，\((y,\delta)\) 与 \(X\) 满足三次约束
\[
4X^{3}-3X^{2}-(4y+2)X+(1+\delta)=0,
\]
其由 \(y=X^2+Y-\tfrac12\) 与 \(\delta=4XY+3X^2-1\) 消去 \(Y\) 得到（见定理 \ref{thm:xi-terminal-zm-delta-spectral-x-inversion} 的证明中的同型消元式）。
将 \(\delta_2\) 视作未知量，与上式及 \(\delta_2=20X^{3}-12X+2+12XY\) 联立，并对 \(X\) 作消元（例如取关于 \(X\) 的结果式）得到所示三次关系；该关系在 \(E\) 上恒成立，故为二阶射流的代数闭式。
\end{proof}

\begin{theorem}[\texorpdfstring{$\delta$}{delta}-闭包的平面五次与五节点归一化]\label{thm:xi-terminal-zm-delta-closure-plane-quintic}
令
\[
F_\delta(T;y):=T^{4}+(8y-3)T^{3}+(2y^{2}-34y-4)T^{2}-y(y-1)\,P_{\mathrm{LY}}(y)\in\ZZ[y][T],
\]
\[
P_{\mathrm{LY}}(y):=256y^{3}+411y^{2}+165y+32.
\]
记仿射平面曲线
\[
\mathcal C_\delta:=\{(y,T)\in\mathbb A^{2}_{(y,T)}:\ F_\delta(T;y)=0\},
\]
并记 \(\overline{\mathcal C}_\delta\subset\mathbb P^{2}\) 为将 \(F_\delta\) 以总次数 \(5\) 齐次化所得的射影闭包。
则成立：
\begin{enumerate}
  \item \(\overline{\mathcal C}_\delta\) 在仿射部分 \(\mathbb A^{2}\) 内恰有五个奇点，并且全为普通双点（节点）；
  \item \(\overline{\mathcal C}_\delta\) 在无穷远处仅有一个点且为光滑点；
  \item \(\overline{\mathcal C}_\delta\) 的归一化为亏格 \(1\) 的光滑射影曲线，并与 \(E\) 在 \(\QQ\) 上同构；
  等价地，
  \[
  \QQ(\overline{\mathcal C}_\delta)\cong \QQ(y,T)/(F_\delta)\cong \QQ(E)=\QQ(X,Y).
  \]
\end{enumerate}
\end{theorem}

\begin{proof}[证明要点]
首先，由结论 \ref{con:xi-terminal-zm-leyang-abel-jacobi-ode} 得 \(\QQ(E)=\QQ(y,\delta)\)，并且 \((y,\delta)\) 满足 \(F_\delta(\delta;y)=0\)；
因此 \(\QQ(y,T)/(F_\delta)\) 的光滑模型为亏格 \(1\) 的曲线，并与 \(E\) 同构（函数域同构确定归一化）。

其次，齐次化取总次数 \(5\)，可写为
\[
F_\delta^{\mathrm{h}}(T,y,Z)=
Z\Bigl(T^{4}+(8y-3Z)T^{3}+(2y^{2}-34yZ-4Z^{2})T^{2}\Bigr)-y(y-Z)P_{\mathrm{LY}}^{\mathrm{h}}(y,Z),
\]
其中
\(
P_{\mathrm{LY}}^{\mathrm{h}}(y,Z)=256y^{3}+411y^{2}Z+165yZ^{2}+32Z^{3}
\)。
令 \(Z=0\) 得 \(F_\delta^{\mathrm{h}}(T,y,0)=-256y^{5}\)，故无穷远处唯一点为 \([T:y:Z]=[1:0:0]\)。
并且
\(
\partial_ZF_\delta^{\mathrm{h}}(1,0,0)=1\neq 0
\)，因此该无穷远点光滑，得到第二条。

对仿射奇点，设
\[
I_{\mathrm{sing}}:=(F_\delta,\partial_TF_\delta,\partial_yF_\delta)\subset\QQ[T,y].
\]
由于奇点必满足 \(F_\delta(T;y)\) 在 \(T\) 上有重根，故其 \(y\)-坐标只能落在判别式零点上。
由结论 \ref{con:xi-terminal-zm-leyang-abel-jacobi-ode} 的判别式分解
\[
\mathrm{Disc}_{T}(F_\delta)=-y(y-1)P_{\mathrm{LY}}(y)\,Q_{5}(y)^{2},
\quad
Q_{5}(y)=4096y^{5}+5376y^{4}-464y^{3}-2749y^{2}-723y+80,
\]
奇点候选仅可能出现在 \(y\in\{0,1\}\cup\{P_{\mathrm{LY}}(y)=0\}\cup\{Q_5(y)=0\}\)。

当 \(T=0\) 时，\(F_\delta(0;y)=-y(y-1)P_{\mathrm{LY}}(y)\)。
直接计算
\[
\partial_yF_\delta(0;y)=-(2y-1)P_{\mathrm{LY}}(y)-y(y-1)P_{\mathrm{LY}}'(y),
\]
并注意 \(P_{\mathrm{LY}}\) 无重根且其根不等于 \(0,1\)，可知在 \(y=0,1\) 与 \(P_{\mathrm{LY}}(y)=0\) 的各点处 \(\partial_yF_\delta(0;y)\neq 0\)；
因此 \((y,T)=(0,0)\)、\((1,0)\) 以及 \(\{P_{\mathrm{LY}}=0\}\times\{0\}\) 仅对应分歧而非平面奇点。
从而一切平面奇点均满足 \(T\neq 0\)，并且其 \(y\)-坐标必满足 \(Q_5(y)=0\)。

在 \(I_{\mathrm{sing}}\) 上可作消元得到
\(
I_{\mathrm{sing}}\cap\QQ[y]=(Q_5(y))
\)，并且在 \(\QQ[T,y]/I_{\mathrm{sing}}\) 中有线性关系
\[
93\,T=P_4(y),
\qquad
P_4(y)=8192y^{4}+5888y^{3}-3680y^{2}-2848y+152.
\]
由于 \(\gcd(Q_5,Q_5')=1\)，故 \(Q_5\) 的五个根互异，从而得到五个互异奇点
\[
\left(y,\ T=\frac{P_4(y)}{93}\right),\qquad Q_5(y)=0.
\]

最后，为判定奇点类型，取二次切锥判别式
\[
\Delta:=(\partial_{Ty}F_\delta)^{2}-(\partial_{TT}F_\delta)(\partial_{yy}F_\delta).
\]
将 \(\Delta\) 在 \(\QQ[T,y]/I_{\mathrm{sing}}\) 中约化得到
\[
\Delta\equiv 2\cdot\bigl(19072y^{4}+51120y^{3}+45129y^{2}+14759y+1032\bigr).
\]
该四次多项式与 \(Q_5\) 互素，故在 \(Q_5(y)=0\) 的各点处 \(\Delta\neq 0\)，从而切锥分解为两条不同直线，五个奇点均为节点。
由于平面五次曲线的算术亏格为 \(\frac{(5-1)(5-2)}{2}=6\)，节点的 \(\delta\)-不变量均为 \(1\)，总和为 \(5\)，故几何亏格为 \(1\)。
与第一段的函数域同构合并，即得第三条。证毕。
\end{proof}

\begin{proposition}[一阶射流的纤维碰撞与 \texorpdfstring{$Q_5$}{Q5} 判据]\label{prop:xi-terminal-zm-leyang-first-jet-fiber-collision-q5}
令 \(\pi_y:E\to\PP^1_y\) 由函数 \(y=X^2+Y-\tfrac12\) 诱导（次数 \(4\)），并在仿射图上定义一阶射流映射
\[
j^{1}y:\ E\longrightarrow \mathbb A^{2}_{(y,T)},
\qquad
P\longmapsto \bigl(y(P),\delta(P)\bigr).
\]
对任意 \(y_0\in\overline{\QQ}\) 满足 \(y_0(y_0-1)P_{\mathrm{LY}}(y_0)\neq 0\)，下列二择一成立：
\begin{enumerate}
  \item 若 \(Q_5(y_0)\neq 0\)，则纤维 \(\pi_y^{-1}(y_0)\) 含四个不同点，且 \(\delta\) 在该纤维上取四个互异值；等价地，\(j^1y\) 在该纤维上为单射。
  \item 若 \(Q_5(y_0)=0\)，则纤维 \(\pi_y^{-1}(y_0)\) 仍含四个不同点，但其中恰有两点 \(P\neq P'\) 满足
  \[
  y(P)=y(P')=y_0,
  \qquad
  \delta(P)=\delta(P')=\frac{8192y_0^{4}+5888y_0^{3}-3680y_0^{2}-2848y_0+152}{93},
  \]
  其余两点在 \(\delta\) 下取值互异。
\end{enumerate}
\end{proposition}
\begin{proof}
固定 \(y=y_0\) 时，\(\delta\) 的可能取值为四次多项式 \(F_\delta(T;y_0)=0\) 的根。
由于 \(y_0(y_0-1)P_{\mathrm{LY}}(y_0)\neq 0\)，由定理 \ref{thm:xi-terminal-zm-leyang-elliptic-structure} 中对 \(\pi_y\) 的分歧刻画可知：\(y_0\) 不是分歧值，故纤维 \(\pi_y^{-1}(y_0)\) 由四个不同点组成。
\par
若 \(Q_5(y_0)\neq 0\)，由定理 \ref{thm:xi-terminal-zm-delta-closure-plane-quintic} 的判别式分解
\(
\Disc_T(F_\delta)=-y(y-1)P_{\mathrm{LY}}(y)\,Q_5(y)^2
\)
知 \(\Disc_T(F_\delta(\cdot;y_0))\neq 0\)，从而 \(F_\delta(\cdot;y_0)\) 有四个互异根；它们分别对应四个片层上的 \(\delta\)-取值，故 \(j^1y\) 在该纤维上单射。
\par
若 \(Q_5(y_0)=0\)，在同一非分歧假设下纤维仍由四个不同点组成；而 \(\Disc_T(F_\delta(\cdot;y_0))=0\) 表明 \(F_\delta(\cdot;y_0)\) 在 \(T\)-轴上出现重根。
由定理 \ref{thm:xi-terminal-zm-delta-closure-plane-quintic} 的消元刻画，该重根满足线性闭式
\(
93\,T=P_4(y)
\)，因此两点在同一 \((y_0,\delta_0)\) 处发生射流识别，且 \(\delta_0=P_4(y_0)/93\)。
其余两点对应的根为单根，故 \(\delta\) 取值互异。
\end{proof}

\begin{proposition}[由 \texorpdfstring{$\delta$}{delta} 乘法诱导的秩 \(4\) 谱结构]\label{prop:xi-terminal-zm-leyang-delta-multiplication-spectral-structure}
令 \(\pi_y:E\to\PP^1_y\) 如上，并令
\[
V:=\pi_{y*}\mathcal O_E.
\]
则 \(V\) 为 \(\PP^1_y\) 上秩 \(4\) 向量丛。
令 \(\Phi:V\to V\) 为“乘以 \(\delta\)”诱导的 \(\mathcal O_{\PP^1_y}\)-线性内自同态。
则在仿射图 \(\mathbb A^1_y\) 上有恒等式
\[
\det\bigl(\eta\cdot\mathrm{Id}_{V}-\Phi\bigr)=F_\delta(\eta;y),
\]
从而 \(F_\delta(\eta;y)=0\) 给出 \((V,\Phi)\) 的谱曲线；其归一化同构于 \(E\)，而平方因子 \(Q_5(y)^2\) 精确记录平面模型中由节点导致的非正则识别。
\end{proposition}
\begin{proof}
在 \(\mathbb A^1_y\) 上，\(V\) 的局部截面可识别为有限平坦代数
\[
\mathcal O_{\mathbb A^1_y}[T]\big/(F_\delta(T;y)),
\]
其中 \(T\) 对应 \(\delta\) 的代数生成元（定理 \ref{thm:xi-terminal-zm-delta-closure-plane-quintic} 的函数域同构）。
在该表示下，\(\Phi\) 恰为“乘以 \(T\)”的乘法算子。
其特征多项式等于 \(T\) 在 \(\QQ(y)\) 上的极小多项式；由于 \(F_\delta\) 为首一四次且不可约，二者一致，从而得到所示行列式恒等式。
谱曲线的归一化与平方因子的几何含义由定理 \ref{thm:xi-terminal-zm-delta-closure-plane-quintic} 与推论 \ref{cor:xi-terminal-zm-delta-discriminant-square-geometry} 给出。
\end{proof}

\begin{theorem}[五节点平面模型的广义雅可比与半阿贝尔分解]\label{thm:xi-terminal-zm-delta-generalized-jacobian-semiabelian}
令 \(\overline{\mathcal C}_\delta\subset\PP^2\) 为定理 \ref{thm:xi-terminal-zm-delta-closure-plane-quintic} 的射影闭包，并令 \(\nu:E\to\overline{\mathcal C}_\delta\) 为归一化映射。
则 \(\overline{\mathcal C}_\delta\) 为 Deligne--Mumford 稳定曲线，算术亏格为 \(6\)，并且其广义雅可比（\(\mathrm{Pic}^0(\overline{\mathcal C}_\delta)\)）为维数 \(6\) 的半阿贝尔簇，满足正合列
\[
1\longrightarrow (\mathbb G_m)^5\longrightarrow \mathrm{Pic}^0(\overline{\mathcal C}_\delta)\longrightarrow E\longrightarrow 0.
\]
更具体地，设五个节点在 \(E\) 上的两两预像为 \(\{P_i,Q_i\}\subset E(\overline{\QQ})\)（\(i=1,\dots,5\)），则该扩张的扩张类由差分除子类 \([P_i-Q_i]\in \mathrm{Pic}^0(E)\cong E\) 决定。
\end{theorem}
\begin{proof}[证明要点]
由定理 \ref{thm:xi-terminal-zm-delta-closure-plane-quintic}，\(\overline{\mathcal C}_\delta\) 仅含五个节点且归一化为椭圆曲线 \(E\)，故
\(
p_a(\overline{\mathcal C}_\delta)=g(E)+5=6
\)，并且稳定性成立。
\par
对节点曲线有标准的单位层短正合列
\[
1\to \mathcal O_{\overline{\mathcal C}_\delta}^{\times}\to \nu_{*}\mathcal O_{E}^{\times}\to \bigoplus_{i=1}^{5}\mathbb G_m\to 1,
\]
其中末项逐节点记录两支粘合的标量比。
取 \(\mathrm H^{1}\) 并取零度部分得到
\[
1\to(\mathbb G_m)^5\to \mathrm{Pic}^0(\overline{\mathcal C}_\delta)\to \mathrm{Pic}^0(E)\to 0,
\]
而 \(\mathrm{Pic}^0(E)\cong E\)，从而得到所示半阿贝尔分解。
扩张类由每个节点的两预像点对 \(\{P_i,Q_i\}\) 的差分除子类决定，这是广义雅可比的标准刻画。
\end{proof}

\begin{proposition}[伴随二次曲线的唯一性与 \texorpdfstring{$Q_5$}{Q5} 的内禀消元刻画]\label{prop:xi-terminal-zm-delta-adjoint-conic-and-elimination}
令
\[
H(T,y):=\frac{1}{T}\,\partial_TF_\delta(T;y)=4T^{2}+(24y-9)T+(4y^{2}-68y-8).
\]
则：
\begin{enumerate}
  \item \(H(T,y)=0\) 在射影意义下为 \(\overline{\mathcal C}_\delta\) 的唯一伴随二次曲线（即通过全部五个节点的二次曲线，按比例唯一）；
  \item 结果式消元满足恒等式
  \[
  \mathrm{Res}_{T}\bigl(F_\delta(T;y),\,H(T,y)\bigr)=Q_{5}(y)^{2}.
  \]
  因而 \(Q_5\) 可被内禀刻画为理想 \((F_\delta,\frac{1}{T}\partial_TF_\delta)\) 的消元因子。
\end{enumerate}
\end{proposition}

\begin{proof}[证明要点]
由于 \(F_\delta\) 不含 \(T\) 的一次项，\(\partial_TF_\delta\) 含因子 \(T\)，故 \(H\) 为二次多项式。
任一节点满足 \(F_\delta=\partial_TF_\delta=\partial_yF_\delta=0\)，并且由定理 \ref{thm:xi-terminal-zm-delta-closure-plane-quintic} 的证明知节点处 \(T\neq 0\)，故节点亦满足 \(H=0\)；
于是 \(H=0\) 通过全部五个节点。

另一方面，次数为 \(5\) 且仅含节点的平面曲线，其正则微分由次数 \(d-3=2\) 的伴随曲线给出；
在亏格 \(1\) 情形，伴随二次曲线空间维数为 \(1\)，故按比例唯一。
因此 \(H=0\) 为唯一伴随二次曲线。

对结果式，使用判别式--结果式恒等式
\[
\mathrm{Disc}_{T}(F_\delta)=\pm\,\mathrm{Res}_{T}(F_\delta,\partial_TF_\delta),
\]
并用 \(\partial_TF_\delta=T\cdot H\) 与结果式乘法性得到
\[
\mathrm{Res}_{T}(F_\delta,\partial_TF_\delta)
=\mathrm{Res}_{T}(F_\delta,T)\,\mathrm{Res}_{T}(F_\delta,H).
\]
注意 \(\mathrm{Res}_{T}(F_\delta,T)=F_\delta(0;y)=-y(y-1)P_{\mathrm{LY}}(y)\)。
与结论 \ref{con:xi-terminal-zm-leyang-abel-jacobi-ode} 的判别式分解合并，即得
\(
\mathrm{Res}_{T}(F_\delta,H)=Q_5(y)^2
\)。
证毕。
\end{proof}

\begin{theorem}[谱坐标 \texorpdfstring{$\lambda=X$}{lambda=X} 的有理反演与极点除子]\label{thm:xi-terminal-zm-delta-spectral-x-inversion}
继续沿用上述记号，并令谱坐标 \(\lambda:=X\)。
在开集
\[
U:=\{(y,T)\in \mathcal C_\delta:\ Q_{5}(y)\neq 0\}
\]
上，有显式有理反演
\[
\lambda=\frac{N(y,T)}{Q_{5}(y)},
\]
其中
\[
\begin{aligned}
N(y,T)=\;&32T^{3}y+19T^{3}+256T^{2}y^{3}+416T^{2}y^{2}+24T^{2}y-81T^{2}\\
&+512Ty^{4}+160Ty^{3}-709Ty^{2}-411Ty-20T\\
&-1792y^{5}-2528y^{4}-914y^{3}-377y^{2}-85y+80.
\end{aligned}
\]
并令
\[
Y:=y-\lambda^{2}+\frac12.
\]
则 \((\lambda,Y)\) 满足椭圆曲线方程
\[
Y^{2}=\lambda^{3}-\lambda+\frac14,
\]
且反演关系
\[
y=\lambda^{2}+Y-\frac12,
\qquad
T=4\lambda Y+3\lambda^{2}-1
\]
成立。
因此在 \(U\) 上，映射
\[
(\lambda,Y)\longmapsto (y,T)=\Bigl(\lambda^{2}+Y-\frac12,\ 4\lambda Y+3\lambda^{2}-1\Bigr)
\]
给出 \(E\dashrightarrow\mathcal C_\delta\) 的双有理参数化，而 \(\lambda=N/Q_{5}\) 给出其在 \(U\) 上的显式逆。
\end{theorem}

\begin{proof}[证明要点]
由结论 \ref{con:xi-terminal-zm-leyang-abel-jacobi-ode} 的消元链条，\(\QQ(E)=\QQ(y,T)\)，故 \(\lambda\in\QQ(y,T)\)。
另一方面，\((\lambda,y)\) 必满足谱四次
\[
\Pi(\lambda,y)=\lambda^{4}-\lambda^{3}-(2y+1)\lambda^{2}+\lambda+y(y+1)=0,
\]
并且由 \(T=4\lambda Y+3\lambda^{2}-1\) 与 \(Y=y-\lambda^{2}+\frac12\) 消去 \(Y\) 得到三次约束
\[
4\lambda^{3}-3\lambda^{2}-(4y+2)\lambda+(1+T)=0.
\]
在 \(\QQ[\lambda,y,T]\) 中对上述两式作消元，可得到线性关系
\[
Q_{5}(y)\lambda-N(y,T)=0,
\]
其中 \(Q_5\) 与 \(N\) 如上。
当 \(Q_5(y)\neq 0\) 时可解得 \(\lambda=N/Q_5\)。

取 \(Y:=y-\lambda^2+\tfrac12\)，则由定义有 \(y=\lambda^2+Y-\tfrac12\)。
将 \(\Pi(\lambda,y)=0\) 代入并整理即得
\(
Y^2=\lambda^3-\lambda+\tfrac14
\)，从而 \((\lambda,Y)\in E\)。
最后将 \((\lambda,Y)\) 代回 \(T=4\lambda Y+3\lambda^2-1\) 得恒等成立，故得双有理互逆。证毕。
\end{proof}

\begin{proposition}[平面射流坐标的显式反演]\label{prop:xi-terminal-zm-delta-jet-plane-explicit-inversion}
\leanverified{paper\_xi\_terminal\_zm\_delta\_jet\_plane\_explicit\_inversion}
在开集
\[
U_{\mathrm{inv}}:=\{(y,T)\in\mathcal C_\delta:\ \mathcal Q(y,T)\neq 0\},
\qquad
\mathcal Q(y,T):=8192y^{4}+5888y^{3}-3680y^{2}-2848y+152-93T,
\]
上，归一化同构 \(\QQ(E)\cong \QQ(y,T)/(F_\delta)\) 的谱坐标可写成显式有理式
\[
X=\frac{\mathcal N(y,T)}{\mathcal Q(y,T)},
\qquad
Y=y-X^{2}+\frac12,
\]
其中
\[
\begin{aligned}
\mathcal N(y,T)=\;&64T^{3}+512T^{2}y^{2}+528T^{2}y-219T^{2}+1024Ty^{3}-288Ty^{2}\\
&-1154Ty-131T-3584y^{4}-2928y^{3}-415y^{2}-929y+152.
\end{aligned}
\]
并且在 \(U_{\mathrm{inv}}\) 上恒等式
\[
y=X^{2}+Y-\frac12,
\qquad
T=4XY+3X^{2}-1
\]
成立。
\end{proposition}
\begin{proof}
由 \(\mathcal Q(y,T)=P_4(y)-93T\) 可知其零点恰对应平面模型中“射流分母退化”的条带；在开集 \(\mathcal Q\neq 0\) 上，可将函数域同构写为 \(X\in\QQ(y,T)\) 的显式有理表达式。
所示 \(\mathcal N/\mathcal Q\) 由对 \((F_\delta(T;y),\,4X^{3}-3X^{2}-(4y+2)X+(1+T))\) 的消元反演给出，并与定理 \ref{thm:xi-terminal-zm-delta-spectral-x-inversion} 的反演在共同定义域上一致。
令 \(Y=y-X^{2}+\tfrac12\) 后，直接代入 \(E:Y^2=X^3-X+\tfrac14\) 与 \(T=4XY+3X^2-1\) 即得两条恒等式。
\end{proof}

\begin{corollary}[判别式分解中平方因子的几何来源]\label{cor:xi-terminal-zm-delta-discriminant-square-geometry}
令 \(\widetilde{\mathcal C}_\delta\to\overline{\mathcal C}_\delta\) 为归一化，\(\pi_y:\widetilde{\mathcal C}_\delta\to\mathbb P^1_y\) 为由坐标 \(y\) 诱导的投影。
则：
\begin{enumerate}
  \item 因子 \(-y(y-1)P_{\mathrm{LY}}(y)\) 恰给出 \(\pi_y\) 的分歧像（有限分歧值），对应于 \(T=0\) 的分支；
  \item 多项式 \(Q_{5}(y)\) 的零点不对应 \(\pi_y\) 的分歧值，而对应于平面闭包 \(\overline{\mathcal C}_\delta\subset\mathbb P^{2}\) 的五个节点在 \(y\)-直线上的投影；
  \item 因而
  \[
  \mathrm{Disc}_{T}(F_\delta(T;y))
  =-\underbrace{y(y-1)P_{\mathrm{LY}}(y)}_{\text{归一化覆盖的分歧贡献}}
  \cdot
  \underbrace{Q_{5}(y)^{2}}_{\text{平面模型节点造成的平方因子}}
  \]
  给出“分歧贡献与平面自交贡献”的正交分解：平方因子完全记录平面化表示中的节点识别，而不改变归一化后的本征分歧集合。
\end{enumerate}
\end{corollary}

\begin{proof}[证明要点]
由命题 \ref{prop:xi-terminal-zm-delta-adjoint-conic-and-elimination} 的结果式分解
\[
\mathrm{Res}_{T}(F_\delta,\partial_TF_\delta)
=\mathrm{Res}_{T}(F_\delta,T)\,\mathrm{Res}_{T}(F_\delta,H),
\]
其中 \(\mathrm{Res}_{T}(F_\delta,T)=F_\delta(0;y)=-y(y-1)P_{\mathrm{LY}}(y)\) 对应 \(T=0\) 处的重根，即 \(\pi_y\) 的分歧像；
而 \(\mathrm{Res}_{T}(F_\delta,H)=Q_5(y)^2\) 对应 \(T\neq 0\) 的约化导数零点。
由定理 \ref{thm:xi-terminal-zm-delta-closure-plane-quintic}，这些点恰为平面模型的节点，其 \(y\)-坐标由 \(Q_5(y)=0\) 给出。
因此平方因子仅记录平面闭包的自交，并不引入新的归一化分歧值。证毕。
\end{proof}

\begin{theorem}[\texorpdfstring{$(5,4)$}{(5,4)} 双齐次闭包的有限节点与无穷远尖点]\label{thm:xi-terminal-zm-delta-bihomogeneous-closure-singularity-spectrum}
令 \(\overline{\mathcal C}_\delta^{\mathrm{bi}}\subset \mathbb P^1_y\times \mathbb P^1_T\) 为 \(F_\delta(T;y)=0\) 的 \((\deg_y,\deg_T)=(5,4)\) 双齐次闭包：取齐次坐标 \((y_0:y_1)\)、\((T_0:T_1)\)，并令
\[
F_\delta^{\mathrm{bi}}(T_0,T_1;y_0,y_1)
:=
y_1^5T_0^4
(8y_0y_1^4-3y_1^5)T_0^3T_1
(2y_0^2y_1^3-34y_0y_1^4-4y_1^5)T_0^2T_1^2
-y_0(y_0-y_1)P_{\mathrm{LY}}^{\mathrm{h}}(y_0,y_1)T_1^4,
\]
其中 \(P_{\mathrm{LY}}^{\mathrm{h}}(y_0,y_1)=256y_0^3+411y_0^2y_1+165y_0y_1^2+32y_1^3\)，并定义
\[
\overline{\mathcal C}_\delta^{\mathrm{bi}}:=\{(y_0:y_1;T_0:T_1):\ F_\delta^{\mathrm{bi}}(T_0,T_1;y_0,y_1)=0\}.
\]
则 \(\overline{\mathcal C}_\delta^{\mathrm{bi}}\) 的奇点谱满足：
\begin{enumerate}
  \item 在仿射图 \(y\neq\infty\)、\(T\neq\infty\) 上，\(\overline{\mathcal C}_\delta^{\mathrm{bi}}\) 恰有五个奇点，且全为普通节点；其 \(y\)-坐标恰为 \(Q_5(y)\) 的五个根，并且对应的 \(T\)-坐标为 \(\delta_\ast(y_0)\)，其中 \(\delta_\ast\) 为定理 \ref{thm:xi-fdelta-subresultant-unique-double-root} 所给出的唯一二重根闭式；
  \item 在边界 \(\{y=\infty\}\cup\{T=\infty\}\) 上，\(\overline{\mathcal C}_\delta^{\mathrm{bi}}\) 仅有一个点，且必为 \((y,T)=(\infty,\infty)\)；在局部坐标 \(z:=1/y\)、\(s:=1/T\) 下，其 Newton 主部与
  \[
  z^{5}-256\,s^{4}=0
  \]
  等价，从而该点为单分支尖点，Puiseux 对为 \((4,5)\)，并且其 \(\delta\)-不变量为 \(6\)；
  \item 除上述五个节点与一个 \((4,5)\)-尖点外，不存在任何额外奇点；
  \item \(\overline{\mathcal C}_\delta^{\mathrm{bi}}\) 的归一化为亏格 \(1\) 的光滑射影曲线，且其函数域与 \(\QQ(E)\) 同构；因而归一化与 \(E\) 在 \(\QQ\) 上同构。
\end{enumerate}
\end{theorem}

\begin{proof}[证明要点]
有限处结论由定理 \ref{thm:xi-terminal-zm-delta-closure-plane-quintic} 与定理 \ref{thm:xi-fdelta-subresultant-unique-double-root} 合并得到：判别式零点给出有限奇点候选，且在 \(Q_5(y)=0\) 处出现唯一二重根 \(\delta_\ast(y_0)\)；同时节点性可由切锥判别式的非退化性判定。

对无穷远处，双齐次方程迫使 \(y=\infty\) 与 \(T=\infty\) 同时发生，故边界候选点唯一。
在局部坐标 \(z=1/y\)、\(s=1/T\) 下，将 \(F_\delta(T;y)=0\) 写成 \(z,s\) 的方程并取最低阶主导平衡，即得 Newton 主部 \(z^5-256s^4\)。
由标准参数化 \(z=t^4\)、\(s=\tfrac{1}{4}t^5\) 可读出 Puiseux 对为 \((4,5)\)，从而 \(\delta\)-不变量为 \((4-1)(5-1)/2=6\)。

最后，\((5,4)\)-型曲线的算术亏格为 \((5-1)(4-1)=12\)；
五个节点贡献 \(\sum\delta_P=5\)，无穷远尖点贡献 \(6\)，故几何亏格为 \(12-(5+6)=1\)。
结合 \(\QQ(\overline{\mathcal C}_\delta^{\mathrm{bi}})=\QQ(y,T)/(F_\delta)\cong \QQ(E)\) 即得归一化同构断言。证毕。
\end{proof}

\begin{proposition}[整闭包导子与判别式的指数平方分解]\label{prop:xi-terminal-zm-delta-conductor-q5-and-discriminant-index-square}
令 \(k:=\QQ(y)\)、\(K:=\QQ(y,T)/(F_\delta)\cong \QQ(E)\)，并设
\[
A:=\QQ[y]\subset K,\qquad
B_0:=A[T]\big/(F_\delta(T;y)).
\]
记 \(B\) 为 \(A\) 在 \(K\) 中的整闭包（等价于 \(\mathcal C_\delta\) 的归一化在 \(y\neq\infty\) 处的坐标环）。
则 \(B_0\subset B\) 的导子理想满足
\[
\mathfrak f:=\operatorname{Ann}_{B}(B/B_0)=(Q_5(y))\subset A,
\]
并且判别式理想满足指数平方公式
\[
\operatorname{Disc}(B_0/A)=\operatorname{Disc}(B/A)\cdot (Q_5(y))^2,
\]
其中 \(\operatorname{Disc}(B_0/A)\) 由 \(F_\delta\) 的 \(T\)-判别式给出：
\[
\operatorname{Disc}(B_0/A)=\bigl(\mathrm{Disc}_T(F_\delta(T;y))\bigr)
=\bigl(-y(y-1)P_{\mathrm{LY}}(y)\,Q_5(y)^2\bigr),
\qquad
\operatorname{Disc}(B/A)=\bigl(-y(y-1)P_{\mathrm{LY}}(y)\bigr).
\]
\end{proposition}

\begin{proof}[证明要点]
由定理 \ref{thm:xi-terminal-zm-delta-closure-plane-quintic}，\(B_0\) 的整闭缺陷恰发生在五个节点之上，且其 \(y\)-坐标为 \(Q_5\) 的五个互异根。
对每个节点的完备局部模型，\(B_0\) 的局部环同构于 \(k[[u,v]]/(uv)\)，其整闭包为 \(k[[u]]\oplus k[[v]]\)，导子由极大理想生成；
在基环 \(A\) 上的像即为 \((y-y_0)\)。
对五点相乘得到 \(\mathfrak f=(Q_5(y))\)。

判别式的指数平方关系 \(\operatorname{Disc}(\text{order})=\operatorname{Disc}(\text{整闭包})\cdot (\text{指数})^2\) 在一维整环的有限扩张中逐点成立；
将五个节点处的局部指数相乘即得 \(Q_5(y)\)，从而得到所示分解。
最后两条显式判别式恒等式分别由
\(\mathrm{Disc}_T(F_\delta)=-y(y-1)P_{\mathrm{LY}}(y)\,Q_5(y)^2\)
以及归一化覆盖的分歧判别得到。证毕。
\end{proof}

\begin{proposition}[\texorpdfstring{$\delta$}{delta}-投影判别式的分歧--结点分离分解]\label{prop:xi-terminal-zm-delta-discy-factorization-a5b5}
将 \(F_\delta(T;y)\) 视为 \(\QQ[T]\) 上关于变量 \(y\) 的五次多项式。则其 \(y\)-判别式在 \(\ZZ[T]\) 中满足精确分解
\begin{equation}\label{eq:xi-terminal-zm-delta-discy-factorization-a5b5}
\boxed{\;
\mathrm{Disc}_{y}\!\bigl(F_\delta(T;y)\bigr)
=16\,(T-4)\,A_5(T)^2\,B_5(T)\; } ,
\end{equation}
其中
\[
\begin{aligned}
A_5(T)
&:=4096T^5-47360T^4+186992T^3-231939T^2-200880T+482112,\\
B_5(T)
&:=50000T^5+136112T^4+60776T^3-26712T^2+38961T+35964.
\end{aligned}
\]
因此 \(\mathrm{Disc}_{y}(F_\delta)\) 的平方因子恰为 \(A_5(T)^2\)，而其平方自由部分为 \(16\,(T-4)\,B_5(T)\)。
\end{proposition}
\begin{proof}
直接计算 \(\mathrm{Disc}_{y}(F_\delta(T;y))\) 并在 \(\ZZ[T]\) 中分解即可得到 \eqref{eq:xi-terminal-zm-delta-discy-factorization-a5b5}。
可复现核验脚本：\texttt{scripts/exp\_fold\_zm\_delta\_projection\_a5b5\_audit.py}。证毕。
\end{proof}

\begin{proposition}[结点子概形的全消元形与坐标域互换同构]\label{prop:xi-terminal-zm-delta-nodes-elimination-a5}
记 \(\mathcal C_\delta\subset\mathbb A^2_{(y,T)}\) 由 \(F_\delta(T;y)=0\) 给出，并令
\[
I_{\mathrm{sing}}:=(F_\delta,\partial_yF_\delta,\partial_TF_\delta)\subset \QQ[y,T].
\]
则 \(\mathcal C_\delta\) 的奇点子概形 \(\mathrm{Sing}(\mathcal C_\delta)\) 由理想 \(I_{\mathrm{sing}}\) 给出，并满足以下两组完全消元形：
\begin{enumerate}
  \item 在消去 \(T\) 后有
  \[
  I_{\mathrm{sing}}\cap \QQ[y]=(Q_5(y)),
  \qquad
  Q_{5}(y)=4096y^{5}+5376y^{4}-464y^{3}-2749y^{2}-723y+80,
  \]
  并且在 \(\QQ[y,T]/I_{\mathrm{sing}}\) 中出现一次关系
  \begin{equation}\label{eq:xi-terminal-zm-delta-node-linear-T}
  \boxed{\;
  93\,T=\bigl(8192y^{4}+5888y^{3}-3680y^{2}-2848y+152\bigr)\; } .
  \end{equation}
  \item 在消去 \(y\) 后有
  \[
  I_{\mathrm{sing}}\cap \QQ[T]=(A_5(T)),
  \]
  其中 \(A_5(T)\) 见命题 \ref{prop:xi-terminal-zm-delta-discy-factorization-a5b5}；并且在 \(\QQ[y,T]/I_{\mathrm{sing}}\) 中出现一次关系
  \begin{equation}\label{eq:xi-terminal-zm-delta-node-linear-y}
  \boxed{\;
  353664\,y+1863680T^{4}-16240384T^{3}+38835472T^{2}+5074731T-77031720=0\; } .
  \end{equation}
\end{enumerate}
从而结点的 \(y\)-坐标全体恰为 \(Q_5\) 的根集，结点的 \(T\)-坐标全体恰为 \(A_5\) 的根集。并且对任一结点 \((y_0,T_0)\) 成立严格域等式
\begin{equation}\label{eq:xi-terminal-zm-delta-node-coordinate-field-swap}
\QQ(y_0)=\QQ(T_0).
\end{equation}
\end{proposition}
\begin{proof}
理想刻画由雅可比判据给出。消元形 \eqref{eq:xi-terminal-zm-delta-node-linear-T} 已在定理 \ref{thm:xi-terminal-zm-delta-closure-plane-quintic} 的证明中出现；
而 \eqref{eq:xi-terminal-zm-delta-node-linear-y} 与 \(I_{\mathrm{sing}}\cap\QQ[T]=(A_5(T))\) 可由对 \(I_{\mathrm{sing}}\) 取 \(y\succ T\) 的字典序 Gröbner 消元得到。
由 \eqref{eq:xi-terminal-zm-delta-node-linear-T} 立得 \(T_0\in\QQ(y_0)\)，由 \eqref{eq:xi-terminal-zm-delta-node-linear-y} 立得 \(y_0\in\QQ(T_0)\)，从而得到 \eqref{eq:xi-terminal-zm-delta-node-coordinate-field-swap}。
可复现核验脚本：\texttt{scripts/exp\_fold\_zm\_delta\_projection\_a5b5\_audit.py}。证毕。
\end{proof}

\begin{theorem}[两个五次多项式的全对称算术单值化与分裂域一致]\label{thm:xi-terminal-zm-delta-a5-b5-galois-s5}
命题 \ref{prop:xi-terminal-zm-delta-discy-factorization-a5b5} 中的 \(A_5(T)\) 与 \(B_5(T)\) 在 \(\QQ\) 上不可约，且其伽罗瓦群均为全对称群 \(S_5\)。
并且 \(Q_5(y)\) 与 \(A_5(T)\) 具有相同分裂域；特别地，
\[
\Gal\bigl(A_5/\QQ\bigr)\cong \Gal\bigl(Q_5/\QQ\bigr)\cong S_5.
\]
\end{theorem}
\begin{proof}
先证 \(A_5\)。取素数 \(p=7\)，由 \(A_5\bmod 7\) 在 \(\FF_7[T]\) 中保持五次不可约可知 \(A_5\) 在 \(\QQ\) 上不可约，故其伽罗瓦群 \(G_A\le S_5\) 传递。
再取素数 \(p=61\)，可验证 \(A_5\bmod 61\) 的分解型为 \((2)(3)\)，故 \(G_A\) 含有一个 \(3\)-循环。
同时 \(\mathrm{Disc}(A_5)\) 在 \(\ZZ\) 中非平方，故 \(G_A\nsubseteq A_5\)。
由含 \(3\)-循环的传递子群分类可得 \(G_A=S_5\)。
\par
对 \(B_5\) 亦同理：取 \(p=59\) 得 \(B_5\bmod 59\) 不可约，故 \(G_B\) 传递；取 \(p=7\) 得分解型 \((2)(3)\) 从而含 \(3\)-循环；并由 \(\mathrm{Disc}(B_5)\) 非平方排除 \(A_5\)，遂 \(G_B=S_5\)。
\par
最后，由命题 \ref{prop:xi-terminal-zm-delta-nodes-elimination-a5} 的域互换同构 \eqref{eq:xi-terminal-zm-delta-node-coordinate-field-swap}：对任一结点 \((y_0,T_0)\) 有 \(\QQ(y_0)=\QQ(T_0)\)。
因此 \(A_5\) 的根生成的分裂域与 \(Q_5\) 的根生成的分裂域一致，从而两者的伽罗瓦群同构。
可复现核验脚本：\texttt{scripts/exp\_fold\_zm\_delta\_projection\_a5b5\_audit.py}。证毕。
\end{proof}

\begin{theorem}[\texorpdfstring{$\delta:E\to\PP^1_\delta$}{delta-map} 的分歧护照与几何单值群]\label{thm:xi-terminal-zm-delta-map-hurwitz-passport-s5}
将 \(\delta=\dd y/\dd u=4XY+3X^2-1\in\QQ(E)\) 视为 \(E\) 上的有理函数，并记由其诱导的有限态射
\[
\pi_\delta:\ E\longrightarrow \PP^1_\delta .
\]
则：
\begin{enumerate}
  \item \(\delta\) 的极点除子满足 \((\delta)_\infty=5O\)。因此 \(\deg(\pi_\delta)=5\)，并且在 \(\delta=\infty\) 处全分歧（分歧指数 \(e_O=5\)）。
  \item \(\pi_\delta\) 的有限分歧点恰由 \(\dd \delta=0\) 给出；其 \(X\)-坐标满足完全分解
  \[
  (X+1)\bigl(100X^5-136X^4+16X^3+40X^2-13X+1\bigr)=0,
  \]
  且其 \(Y\)-坐标由
  \[
  Y=\frac{-10X^3+6X-1}{6X}
  \]
  唯一确定。上述六点均为简单分歧点（分歧指数 \(2\)）。
  \item 有限分歧值集合为
  \[
  \{\delta=4\}\ \cup\ \{\delta\in\overline{\QQ}:\ B_5(\delta)=0\},
  \]
  其中 \(B_5\) 如命题 \ref{prop:xi-terminal-zm-delta-discy-factorization-a5b5}。
  更精确地，消去 \(X\) 得到严格结果式恒等式
  \begin{equation}\label{eq:xi-terminal-zm-delta-branch-resultant-b5}
  \boxed{\;
  \mathrm{Res}_X\!\left(
  100X^5-136X^4+16X^3+40X^2-13X+1,\ \ 3\delta+20X^3-9X^2-12X+5
  \right)=4860\,B_5(\delta)\; }.
  \end{equation}
  \item 该覆盖的几何单值群同构于 \(S_5\)。等价地，函数域扩张 \(\QQ(E)/\QQ(\delta)\) 的伽罗瓦闭包群为 \(S_5\)。
\end{enumerate}
\end{theorem}
\begin{proof}
由 \(\operatorname{ord}_O(X)=-2\)、\(\operatorname{ord}_O(Y)=-3\) 得
\(
\operatorname{ord}_O(\delta)=-(2+3)=-5
\)，从而 \((\delta)_\infty=5O\) 并得到第一条。
\par
取 \(\omega=\dd X/(2Y)\) 并用 \(\dd Y=(3X^2-1)\omega\)、\(\dd X=2Y\omega\) 微分 \(\delta=4XY+3X^2-1\)，得
\[
\dd\delta=2\bigl(10X^3+6XY-6X+1\bigr)\,\omega.
\]
从而有限分歧点由 \(10X^3+6XY-6X+1=0\) 给出，并解得所示 \(Y\) 的闭式；
代回 \(Y^2=X^3-X+\tfrac14\) 化简即得 \(X\)-坐标多项式的完全分解，得到第二条。
由 Riemann--Hurwitz 公式，度为 \(5\) 的覆盖 \(E\to\PP^1\) 的总分歧度为 \(2\deg(\pi_\delta)=10\)；
无穷远点贡献 \(e_O-1=4\)，故有限分歧总贡献为 \(6\)，与上述六个有限分歧点一致，从而它们均为简单分歧点。
\par
第三条中，\(\delta\) 在 \(X=-1\) 的分歧点处取值为 \(4\)；
对五次因子根消去 \(X\) 即得 \eqref{eq:xi-terminal-zm-delta-branch-resultant-b5}，从而其分歧值满足 \(B_5(\delta)=0\)。
\par
第四条由分歧循环型给出：\(\delta=\infty\) 处全分歧对应一个 \(5\)-循环；任一有限分歧值处为简单分歧，对应换位。
包含 \(5\)-循环与换位的传递子群必为 \(S_5\)。证毕。
\end{proof}

\begin{proposition}[特征 \texorpdfstring{$37$}{37} 上的分歧塌缩：\texorpdfstring{$B_5$}{B5} 的双重根]\label{prop:xi-terminal-zm-delta-badprime-37-b5-double-root}
\leanverified{paper\_xi\_terminal\_zm\_delta\_badprime\_37\_b5\_double\_root}
令 \(B_5(\delta)\) 如命题 \ref{prop:xi-terminal-zm-delta-discy-factorization-a5b5}。
则在 \(\FF_{37}[\delta]\) 中有精确分解
\[
B_5(\delta)\equiv 13\,\delta^2\bigl(\delta^3+2\delta^2-4\delta+3\bigr)\pmod{37}.
\]
因此 \(\pi_\delta\) 的有限分歧值在特征 \(37\) 的约化中发生并合，等价地，\(\delta\)-分歧多项式在 \(p=37\) 处出现重根并导致该覆盖不具良分歧约化。
\end{proposition}
\begin{proof}
对 \(B_5\) 的系数作模 \(37\) 化简并因式分解即可得到所示恒等式。
可复现核验脚本：\texttt{scripts/exp\_fold\_zm\_delta\_projection\_a5b5\_audit.py}。证毕。
\end{proof}

\begin{corollary}[有限节点集的伽罗瓦传递性]\label{cor:xi-terminal-zm-delta-nodes-single-galois-orbit}
\(\overline{\mathcal C}_\delta^{\mathrm{bi}}\) 的五个有限节点在 \(\operatorname{Gal}(\overline{\QQ}/\QQ)\) 作用下构成单一轨道；特别地，\(\overline{\mathcal C}_\delta^{\mathrm{bi}}\) 在 \(\QQ\) 上不存在有限有理奇点。
\end{corollary}

\begin{proof}
由定理 \ref{thm:xi-terminal-zm-delta-bihomogeneous-closure-singularity-spectrum}，有限节点的 \(y\)-坐标恰为 \(Q_5\) 的根；
由定理 \ref{thm:xi-q5-galois-s5}，\(\operatorname{Gal}(Q_5/\QQ)\cong S_5\) 在根集上传递，故节点集亦为单一伽罗瓦轨道。
若存在 \(\QQ\)-有理有限节点，则其 \(y\)-坐标给出 \(Q_5\) 的有理根，与不可约性矛盾。证毕。
\end{proof}

上述五个有限节点在归一化曲线 \(E\) 上各分裂为两支预像，从而在 \(E(\overline{\QQ})\) 上给出次数 \(10\) 的闭点集合。
其谱坐标 \(X\)-层不仅可被一次消元化为显式十次式，而且该十次层的总体伽罗瓦闭包达到与五节点块系相容的最大对称。

\begin{theorem}[伴随二次曲线与 \texorpdfstring{$\overline{\mathcal C}_\delta$}{Cdelta} 的纯节点交]\label{thm:xi-terminal-zm-delta-adjoint-conic-pure-nodal-intersection}
令 \(H(T,y)\) 为命题 \ref{prop:xi-terminal-zm-delta-adjoint-conic-and-elimination} 的伴随二次曲线（按比例唯一），并视其为射影二次曲线 \(\{H=0\}\subset\PP^2\)。
则在射影平面中
\[
\overline{\mathcal C}_\delta\cap\{H=0\}
\]
由 \(\overline{\mathcal C}_\delta\) 的五个有限节点组成，且每个节点的交点重数恰为 \(2\)；特别地，除节点外二次曲线 \(\{H=0\}\) 与 \(\overline{\mathcal C}_\delta\) 无其他交点。
\end{theorem}
\begin{proof}
由命题 \ref{prop:xi-terminal-zm-delta-adjoint-conic-and-elimination}，\(\{H=0\}\) 通过 \(\overline{\mathcal C}_\delta\) 的全部五个有限节点。
对任一有限节点 \(P\in\overline{\mathcal C}_\delta\)，其解析邻域中 \(\overline{\mathcal C}_\delta\) 由两条互异分支组成；任意通过 \(P\) 的曲线与两分支的交点重数分别至少为 \(1\)，故
\(
I_P(\overline{\mathcal C}_\delta,H)\ge 2
\)。
从而
\[
\sum_{P\in \overline{\mathcal C}_\delta\cap\{H=0\}} I_P(\overline{\mathcal C}_\delta,H)\ \ge\ 5\times 2\ =\ 10.
\]
另一方面，\(\deg(\overline{\mathcal C}_\delta)=5\) 且 \(\deg(H)=2\)，由 Bézout 定理总交数（计重数）等于 \(10\)，故上式取等号。
于是交点仅为这五个节点且每点交点重数恰为 \(2\)。证毕。
\end{proof}

\begin{theorem}[节点预像十点在椭圆曲线群上的严格求和恒等式]\label{thm:xi-terminal-zm-delta-node-preimage-sum-zero}
令 \(\nu:E\to\overline{\mathcal C}_\delta\) 为归一化映射。
对五个有限节点 \(P_1,\dots,P_5\in\overline{\mathcal C}_\delta\) 及其在 \(E\) 上的两点预像
\[
\nu^{-1}(P_i)=\{P_i',P_i''\}\subset E(\overline{\QQ}),
\]
则在椭圆曲线群律下成立恒等式
\[
\sum_{i=1}^{5}\bigl(P_i'+P_i''\bigr)=O,
\]
其中 \(O\) 为 \(E\) 的单位元。
\end{theorem}
\begin{proof}
令 \(h:=H(T,y)|_{E}\in\QQ(E)^\times\)，其中 \(H\) 如命题 \ref{prop:xi-terminal-zm-delta-adjoint-conic-and-elimination}，而 \(T=\delta=4XY+3X^2-1\)、\(y=X^2+Y-\tfrac12\) 为 \(E\) 上的有理函数（定理 \ref{thm:xi-terminal-zm-delta-closure-plane-quintic} 的记号）。
由定理 \ref{thm:xi-terminal-zm-delta-adjoint-conic-pure-nodal-intersection}，\(\{H=0\}\cap \overline{\mathcal C}_\delta\) 仅由五个节点组成且每点交点重数为 \(2\)。
拉回到归一化 \(E\) 得到 \(h\) 的零点除子为
\[
(h)_0=\sum_{i=1}^{5}(P_i'+P_i'').
\]
另一方面，由 \(\operatorname{ord}_O(X)=-2\)、\(\operatorname{ord}_O(Y)=-3\) 得 \(\operatorname{ord}_O(\delta)=-5\)，故 \(H\) 的主导项 \(4T^2\) 给出
\[
(h)_\infty=10\,O.
\]
从而
\[
(h)=\sum_{i=1}^{5}(P_i'+P_i'')-10\,O.
\]
椭圆曲线上的主除子在 \(E(\overline{\QQ})\) 的群律下满足“零点之和等于极点之和”；由于 \(10\,O=O\)，遂得 \(\sum_{i=1}^5(P_i'+P_i'')=O\)。证毕。
\end{proof}

\begin{theorem}[节点预像的二次中介式与十次消元]\label{thm:xi-terminal-zm-delta-node-preimage-elimination-r10}
令 \(\nu:E\to\overline{\mathcal C}_\delta\) 为定理 \ref{thm:xi-terminal-zm-delta-generalized-jacobian-semiabelian} 中的归一化映射，并令
\[
\Sigma:=\nu^{-1}\!\bigl(\mathrm{Sing}(\overline{\mathcal C}_\delta)\bigr)\subset E(\overline{\QQ})
\]
为有限节点的规范化预像（按重数计为 \(10\) 个点）。
则在 \((X,y)\)-平面上成立以下完全消元刻画：
\begin{enumerate}
  \item 记 \(Q_5(y)\) 为节点消元五次式
  \[
  Q_{5}(y)=4096y^{5}+5376y^{4}-464y^{3}-2749y^{2}-723y+80.
  \]
  则 \(\Sigma\) 的 \((X,y)\)-投影恰由 \(Q_5(y)=0\) 与下式联立刻画：
  \begin{equation}\label{eq:xi-terminal-zm-delta-node-preimage-quadratic-mediator}
  \boxed{\;
  279X^{2}+\bigl(65536y^{4}+47104y^{3}-35392y^{2}-22412y+1774\bigr)X
  +\bigl(110592y^{4}+67584y^{3}-60096y^{2}-32496y+3075\bigr)=0\; } .
  \end{equation}
  \item 进一步对 \(y\) 消元后，\(\Sigma\) 的 \(X\)-坐标全体为十次多项式 \(R_{10}(X)\) 的全部根：
  \begin{equation}\label{eq:xi-terminal-zm-delta-node-preimage-r10}
  \boxed{\;
  \begin{aligned}
  R_{10}(X)=\;&4096X^{10}-10240X^{9}+4864X^{8}+16384X^{7}-26320X^{6}\\
  &+2744X^{5}+20657X^{4}-14144X^{3}-4730X^{2}+5768X+505.
  \end{aligned}\; } .
  \end{equation}
  等价地，
  \[
  X(\Sigma)=\{\alpha\in\overline{\QQ}:\ R_{10}(\alpha)=0\},
  \]
  按重数计为 \(10\) 个点。
\end{enumerate}
\leanverified{paper\_xi\_terminal\_zm\_delta\_node\_preimage\_elimination\_r10}
\end{theorem}

\begin{proof}[证明要点]
由 \(y=X^{2}+Y-\tfrac12\) 消去 \(Y\) 并代入 \(E:Y^{2}=X^{3}-X+\tfrac14\) 得到谱四次消元式
\[
\Pi(X,y)=X^{4}-X^{3}-(2y+1)X^{2}+X+y(y+1)=0,
\]
其即 \(\pi_y:E\to\PP^1_y\) 的代数消去式（见定理 \ref{thm:xi-terminal-zm-delta-spectral-x-inversion}）。
\par
另一方面，由命题 \ref{prop:xi-terminal-zm-delta-nodes-elimination-a5} 的节点消元形 \eqref{eq:xi-terminal-zm-delta-node-linear-T} 可知：有限节点满足 \(Q_5(y)=0\) 且其 \(\delta\)-坐标满足一次闭式 \(93\,\delta=P_4(y)\)，其中
\[
P_4(y)=8192y^{4}+5888y^{3}-3680y^{2}-2848y+152.
\]
将 \(\delta=4XY+3X^{2}-1\) 与 \(Y=y-X^{2}+\tfrac12\) 代入上式，把节点条件拉回到 \(\QQ[X,y]\) 上；再与 \(\Pi(X,y)=0\) 联立形成有限生成理想。
对该理想作字典序 Gröbner 消元（例如取 \(X\succ y\)）即可得到消元基底在 \(\QQ[X,y]\) 中由 \(\{Q_5(y),\ \eqref{eq:xi-terminal-zm-delta-node-preimage-quadratic-mediator}\}\) 生成，从而得到第一条。
\par
再对 \(y\) 消元（例如对 \(Q_5\) 与 \eqref{eq:xi-terminal-zm-delta-node-preimage-quadratic-mediator} 取关于 \(y\) 的结果式并约去整系数公因子）即得十次式 \eqref{eq:xi-terminal-zm-delta-node-preimage-r10}。
可复现核验脚本：\texttt{scripts/exp\_xi\_leyang\_node\_preimage\_r10\_wreath\_audit.py}。证毕。
\end{proof}

\begin{remark}[节点层的最小代数连结]\label{rem:xi-terminal-zm-delta-node-preimage-minimal-algebraic-bridge}
五次层 \(Q_5(y)=0\) 与十次层 \(R_{10}(X)=0\) 之间通过中介二次式 \eqref{eq:xi-terminal-zm-delta-node-preimage-quadratic-mediator} 严格连接；
该二次式在每个节点 \(y_0\) 之上给出两支预像的 \(X\)-坐标，从而把节点分裂方向转写为 \((C_2)^5\) 型的二次层数据，并为数域计算与分支单值分析提供可直接使用的最小代数表示。
\end{remark}

\begin{theorem}[节点纤维的二次子结式方程与八次判别式]\label{thm:xi-terminal-zm-delta-node-fiber-subresultant-d8}
\leanverified{paper\_xi\_terminal\_zm\_delta\_node\_fiber\_subresultant\_d8}
对节点坐标 \(y_0\)（满足 \(Q_5(y_0)=0\)），节点两预像点在 \(E\) 上的 \(X\)-坐标恰为关于 \(X\) 的二次方程
\[
G_{\mathrm{node}}(X;y_0)=0
\]
的两根，其中
\[
\begin{aligned}
G_{\mathrm{node}}(X;y):={}&(1488y+1023)X^2\\
&+(32768y^4+23552y^3-14720y^2-11020y-322)X\\
&-(8192y^4+5888y^3-2192y^2-1360y+245).
\end{aligned}
\]
其判别式满足严格恒等式
\[
\Disc_X\bigl(G_{\mathrm{node}}(X;y)\bigr)=16\,D_8(y),
\]
其中
\[
\begin{aligned}
D_8(y)={}&67108864 y^8+96468992 y^7-25624576 y^6-85426176 y^5\\
&-15933952 y^4+20019264 y^3+7115981 y^2+186875 y+69139.
\end{aligned}
\]
并且 \(D_8\) 与 \(Q_5\) 在 \(\QQ[y]\) 中互素；从而对任意节点坐标 \(y_0\) 有 \(D_8(y_0)\neq 0\)，故两根互异且对应两点预像不同。
\end{theorem}
\begin{proof}[证明要点]
在 \((X,y)\)-平面上，点 \((X,Y)\in E\) 与 \((y,T)\in \overline{\mathcal C}_\delta\) 的关系满足
\[
\Pi(X,y)=X^{4}-X^{3}-(2y+1)X^{2}+X+y(y+1)=0
\]
以及
\[
4X^{3}-3X^{2}-(4y+2)X+(1+T)=0
\]
（分别由 \(y=X^2+Y-\tfrac12\) 与 \(T=4XY+3X^2-1\) 消去 \(Y\) 得）。
在节点处另有一次闭式 \(93\,T=P_4(y)\)（命题 \ref{prop:xi-terminal-zm-delta-nodes-elimination-a5}）。
因此节点预像的 \(X\)-坐标恰为两多项式
\[
\Pi(X,y)\quad\text{与}\quad 4X^{3}-3X^{2}-(4y+2)X+\Bigl(1+\tfrac{P_4(y)}{93}\Bigr)
\]
在 \(\QQ(y)[X]\) 中的公共根。
对上述二式作关于 \(X\) 的子结式计算可得：当且仅当 \(Q_5(y)=0\) 时，一次子结式整体消失，而二次子结式在约去公因子后恰为所示 \(G_{\mathrm{node}}(X;y)\)；
于是对节点坐标 \(y_0\)，公共因子在 \(X\) 中的次数为 \(2\)，其根即为两点预像的 \(X\)-坐标。
\par
判别式恒等式 \(\Disc_X(G_{\mathrm{node}})=16D_8\) 为直接展开化简所得；而 \(\gcd(D_8,Q_5)=1\) 可由欧几里得算法验证。
可复现核验脚本：\texttt{scripts/exp\_xi\_leyang\_node\_preimage\_r10\_wreath\_audit.py}。证毕。
\end{proof}

\begin{proposition}[节点切锥判别式的范数非平方与分裂方向的二次不可平凡性]\label{prop:xi-terminal-zm-delta-node-tangent-discriminant-norm-nonsquare}
沿用定理 \ref{thm:xi-terminal-zm-delta-closure-plane-quintic} 的平面模型 \(\mathcal C_\delta\subset\mathbb A^2_{(y,T)}\) 与节点理想
\[
I_{\mathrm{sing}}=(F_\delta,\partial_yF_\delta,\partial_TF_\delta)\subset\QQ[y,T],
\qquad
I_{\mathrm{sing}}\cap\QQ[y]=(Q_5(y)).
\]
对任一有限节点 \(P_0=(y_0,T_0)\)（其中 \(Q_5(y_0)=0\) 且 \(T_0\neq 0\)），记其坐标域
\[
\kappa_0:=\QQ(y_0)=\QQ(T_0)
\]
（命题 \ref{prop:xi-terminal-zm-delta-nodes-elimination-a5}），并定义切锥二次型的判别式平方类
\[
\Delta(P_0):=\bigl((\partial_{Ty}F_\delta)^2-(\partial_{TT}F_\delta)(\partial_{yy}F_\delta)\bigr)(y_0,T_0)\in \kappa_0^\times/(\kappa_0^\times)^2.
\]
则 \(\Delta(P_0)\) 由 \(y_0\) 的四次多项式值给出：存在
\begin{equation}\label{eq:xi-terminal-zm-delta-node-tangent-discriminant-Dy}
D(y):=2\bigl(19072y^4+51120y^3+45129y^2+14759y+1032\bigr)\in\QQ[y]
\end{equation}
使得
\[
\Delta(P_0)=D(y_0)\in\kappa_0^\times/(\kappa_0^\times)^2.
\]
并且对每个节点 \(P_0\)，元素 \(D(y_0)\in\kappa_0^\times\) 不是平方。
更精确地，其范数满足显式恒等式
\begin{equation}\label{eq:xi-terminal-zm-delta-node-tangent-discriminant-norm}
\mathrm N_{\kappa_0/\QQ}\!\bigl(D(y_0)\bigr)
=\frac{\Res_y(Q_5(y),D(y))}{4096^4}
=-\frac{3^{17}\cdot 11^{4}\cdot 31^{2}\cdot 727^{2}}{2^{16}}\in\QQ^\times,
\end{equation}
从而归一化映射 \(\nu:\widetilde{\mathcal C}_\delta\to\mathcal C_\delta\) 在 \(P_0\) 处的两条分支方向只能在二次扩张
\[
\kappa_0\bigl(\sqrt{D(y_0)}\bigr)/\kappa_0
\]
上被区分；等价地，两条分支点在 \(\kappa_0\) 上不可分别定义，而被该二次扩张的非平凡自同构交换。
\end{proposition}
\begin{proof}
定理 \ref{thm:xi-terminal-zm-delta-closure-plane-quintic} 的证明已给出：在商环 \(\QQ[y,T]/I_{\mathrm{sing}}\) 中，
\[
(\partial_{Ty}F_\delta)^2-(\partial_{TT}F_\delta)(\partial_{yy}F_\delta)\equiv D(y),
\]
其中 \(D(y)\) 如 \eqref{eq:xi-terminal-zm-delta-node-tangent-discriminant-Dy}，故 \(\Delta(P_0)=D(y_0)\)。
\par
范数恒等式 \eqref{eq:xi-terminal-zm-delta-node-tangent-discriminant-norm} 由一般的范数--结式公式给出：对 \(Q_5\) 的任一根 \(y_0\)，
\[
\mathrm N_{\kappa_0/\QQ}\!\bigl(D(y_0)\bigr)=\frac{\Res_y(Q_5(y),D(y))}{\mathrm{LC}(Q_5)^{\deg D}}
=\frac{\Res_y(Q_5(y),D(y))}{4096^4}.
\]
直接计算结果式并分解得到右端显式值。
该有理数为负且在素数 \(3\) 处具有奇数赋值，故不为 \(\QQ^{\times 2}\)；
若 \(D(y_0)\) 在 \(\kappa_0\) 中为平方，则其范数必在 \(\QQ\) 中为平方，矛盾。
因此 \(D(y_0)\notin(\kappa_0^\times)^2\)。
\par
节点的局部切锥由 \(\Delta(P_0)\) 控制：当且仅当 \(D(y_0)\) 为平方时，切锥两条切线在 \(\kappa_0\) 上分裂为两条可分别定义的直线；
在非平方情形，二者仅在二次扩张 \(\kappa_0(\sqrt{D(y_0)})\) 上分裂，且被 Galois 自同构交换。
这等价于归一化纤维两分支方向的最小分裂域为 \(\kappa_0(\sqrt{D(y_0)})\)。
\end{proof}

\begin{theorem}[节点分裂方向的花环积伽罗瓦群与十次结果式模型]\label{thm:xi-terminal-zm-delta-node-tangent-wreath}
令 \(K\) 为 \(Q_5(y)\) 的分裂域，并取其根为 \(y_1,\dots,y_5\)。
置
\[
\Delta_i:=D(y_i)\in K^\times,
\qquad
L:=K\bigl(\sqrt{\Delta_1},\dots,\sqrt{\Delta_5}\bigr).
\]
则 \(L/\QQ\) 为 Galois 扩张，且
\[
\Gal(L/\QQ)\cong C_2\wr S_5\cong W(B_5),
\qquad
|\Gal(L/\QQ)|=2^5\cdot 5!=3840.
\]
此外，\(L\) 是十次多项式
\begin{equation}\label{eq:xi-terminal-zm-delta-node-tangent-resultant-P10}
P(z):=\Res_y\!\bigl(Q_5(y),\,z^2-D(y)\bigr)\in\ZZ[z]
\end{equation}
的分裂域；其全部根为 \(\{\pm\sqrt{\Delta_i}\}_{i=1}^5\)。
并且 \(P\) 具有仅含偶次幂的显式展开：
\[
\begin{aligned}
P(z)=\;&16777216\,z^{10}-1897083633664\,z^{8}+7820079733545728\,z^{6}\\
&\quad+8189770598748183555\,z^{4}+2512279077174710784864\,z^{2}
+245846622126647065658112.
\end{aligned}
\]
\end{theorem}
\begin{proof}[证明要点]
由 \(\Gal(K/\QQ)\cong S_5\)（定理 \ref{thm:xi-q5-galois-s5}）可知 \(S_5\) 传递置换 \(\{y_i\}\)，从而传递置换 \(\{\Delta_i\}\) 与二次因子 \(\{z^2-\Delta_i\}\)。
因此 \(L/\QQ\) 的伽罗瓦群自然嵌入带符号置换群 \((C_2)^5\rtimes S_5\cong C_2\wr S_5\)。
\par
令 \(V:=K^\times/(K^\times)^2\) 并置
\(
W:=\langle[\Delta_1],\dots,[\Delta_5]\rangle_{\FF_2}\subset V
\)。
则 \(\Gal(L/K)\cong W^\vee\)，从而 \([L:K]=2^{\dim_{\FF_2}W}\)。
由命题 \ref{prop:xi-terminal-zm-delta-node-tangent-discriminant-norm-nonsquare}，
\[
\prod_{i=1}^{5}\Delta_i
=\mathrm N_{\QQ(y_1)/\QQ}\!\bigl(D(y_1)\bigr)\in\QQ^\times/(\QQ^\times)^2
\]
给出 \(W\) 的非零 \(S_5\)-不变量，因此 \(W^{S_5}\neq 0\)。
同时，可取一个在 \(K\) 中完全分裂的素数处检验剩余平方性不一致，从而排除 \(W\) 为一维平凡子模；
在 \(\FF_2\) 的置换表示中，这迫使 \(W\) 等于五维置换模本身，故 \(\dim_{\FF_2}W=5\)，即 \([L:K]=2^5\)，得到
\(\Gal(L/\QQ)\cong (C_2)^5\rtimes S_5\)。
\par
最后，\eqref{eq:xi-terminal-zm-delta-node-tangent-resultant-P10} 的根由结果式定义立即等于 \(\{\pm\sqrt{\Delta_i}\}\)；
因此其分裂域恰为 \(L\)。
\end{proof}

\begin{theorem}[十次分支多项式的带符号循环分解律与密度刻画]\label{thm:xi-terminal-zm-delta-node-tangent-signed-cycle-factorization}
沿用定理 \ref{thm:xi-terminal-zm-delta-node-tangent-wreath} 的记号，并记
\(
G=\Gal(L/\QQ)\cong (C_2)^5\rtimes S_5
\)。
对任意在 \(L/\QQ\) 中不分歧的素数 \(p\)，以标准实现将 \(\mathrm{Frob}_p\) 视为带符号置换 \((\varepsilon,\sigma)\in (C_2)^5\rtimes S_5\)。
将 \(\sigma\) 分解为不交循环 \(c=(i_1\ i_2\ \cdots\ i_\ell)\)，并定义循环符号
\[
s(c):=\varepsilon_{i_1}\varepsilon_{i_2}\cdots \varepsilon_{i_\ell}\in\{\pm 1\}.
\]
则 \(P(z)\bmod p\) 在 \(\FF_p[z]\) 中的不可约因子次数完全由循环数据 \(\{(\ell,s(c))\}\) 决定：
\begin{itemize}
  \item 若 \(c\) 的长度为 \(\ell\) 且 \(s(c)=+1\)，则贡献两个次数为 \(\ell\) 的不可约因子；
  \item 若 \(c\) 的长度为 \(\ell\) 且 \(s(c)=-1\)，则贡献一个次数为 \(2\ell\) 的不可约因子。
\end{itemize}
特别地，当 \(\sigma\) 为 \(5\)-循环时，\(P(z)\bmod p\) 或分解为两个五次不可约因子，或保持十次不可约；且在全部不分歧素数中这两种情形的自然密度均为 \(1/10\)。
\end{theorem}
\begin{proof}
群 \(G\) 在根集合 \(\{\pm\sqrt{\Delta_i}\}_{i=1}^5\) 上的作用为
\(
\sqrt{\Delta_i}\mapsto \varepsilon_i\sqrt{\Delta_{\sigma(i)}}
\)。
沿循环 \(c=(i_1\cdots i_\ell)\) 迭代 \(\ell\) 次得到
\(
\sqrt{\Delta_{i_1}}\mapsto s(c)\sqrt{\Delta_{i_1}}
\)。
因此当 \(s(c)=+1\) 时，\(\sqrt{\Delta_{i_1}}\) 的轨道长度为 \(\ell\)，而 \(-\sqrt{\Delta_{i_1}}\) 形成另一独立轨道；对应于两个次数为 \(\ell\) 的不可约因子。
当 \(s(c)=-1\) 时，\(\sqrt{\Delta_{i_1}}\) 经 \(\ell\) 次变为其相反数，需 \(2\ell\) 次才回到自身，轨道长度为 \(2\ell\)，对应于一个次数为 \(2\ell\) 的不可约因子。
对 \(\sigma\) 的各不交循环相乘得到一般分解型结论。
\par
在 \(\sigma\) 为 \(5\)-循环时，仅有一个循环，故只剩两种分解情形。
计数密度：\(S_5\) 中 \(5\)-循环个数为 \(24\)，对给定 \(5\)-循环，\(\varepsilon\in\{\pm1\}^5\) 有 \(32\) 种，其中 \(\prod_{i=1}^5\varepsilon_i\) 取 \(\pm1\) 各 \(16\) 种；故两类带符号 \(5\)-循环各有 \(24\cdot 16=384\) 个群元素，占 \(|G|=3840\) 的比例均为 \(384/3840=1/10\)。
由 Chebotarev 密度定理得到对应素数密度。证毕。
\end{proof}

\begin{conclusion}[节点分裂方向域的二次子域全分类]\label{con:xi-terminal-zm-delta-node-tangent-quadratic-subfields}
沿用定理 \ref{thm:xi-terminal-zm-delta-node-tangent-wreath} 的记号，并令 \(L/\QQ\) 为节点分裂方向域。
则 \(L\) 的二次子域恰有三个，且可显式取为
\[
\QQ(\sqrt{-67611}),\qquad
\QQ(\sqrt{-3}),\qquad
\QQ(\sqrt{22537})\ \ (22537=31\cdot 727).
\]
\end{conclusion}
\begin{proof}
由定理 \ref{thm:xi-terminal-zm-delta-node-tangent-wreath}，\(\Gal(L/\QQ)\cong (C_2)^5\rtimes S_5\)。
其交换化满足
\[
\bigl((C_2)^5\rtimes S_5\bigr)^{\mathrm{ab}}\cong C_2\times C_2,
\]
其中一个 \(C_2\) 来自 \(S_5\) 的符号表示，另一个 \(C_2\) 来自 \((C_2)^5\) 的“总翻转”商（置换交换子将五个坐标差消去，仅保留全体同时取反的类）。
因此 \(\mathrm{Hom}(\Gal(L/\QQ),C_2)\) 恰有三个非零元，对应恰有三个二次子域。
\par
对其显式生成元：
\begin{enumerate}
  \item 由投影到 \(S_5\) 后取符号表示得到的二次特征，其不动域为 \(K^{A_5}\)（其中 \(K\) 为 \(Q_5\) 的分裂域）。
  因为 \(\Gal(K/\QQ)\cong S_5\)（定理 \ref{thm:xi-q5-galois-s5}），该二次子域由 \(\sqrt{\Disc(Q_5)}\) 生成，从而等于 \(\QQ(\sqrt{-67611})\)（定理 \ref{thm:xi-leyang-ply-q5-linearly-disjoint-chebotarev} 的证明）。
  \item 由 \((C_2)^5\) 上“总翻转”得到的二次特征，其不动域由 \(\sqrt{\Delta_1\cdots \Delta_5}\) 生成。
  由命题 \ref{prop:xi-terminal-zm-delta-node-tangent-discriminant-norm-nonsquare} 的范数恒等式 \eqref{eq:xi-terminal-zm-delta-node-tangent-discriminant-norm}，
  \[
  \Delta_1\cdots \Delta_5
  =\mathrm N_{\QQ(y_1)/\QQ}\!\bigl(D(y_1)\bigr)\in -3\cdot(\QQ^\times)^2,
  \]
  故该二次子域等于 \(\QQ(\sqrt{-3})\)。
  \item 第三条二次特征为上述两条的乘积；其不动域为二次合成域中的第三个二次子域，即
  \[
  \QQ\!\bigl(\sqrt{(-67611)\cdot(-3)}\bigr)=\QQ(\sqrt{202833})=\QQ(\sqrt{22537}).
  \]
\end{enumerate}
证毕。
\end{proof}

\begin{remark}[完全分裂素数处的剩余平方性模式]\label{rem:xi-terminal-zm-delta-node-tangent-split-prime-pattern}
存在无穷多素数 \(p\) 使 \(Q_5\) 在 \(\FF_p\) 上完全分裂（Chebotarev 定理）。
在此类素数中，可出现节点切锥判别式 \(D(y_i)\) 的剩余平方性不一致的情形：例如在 \(p=3229\) 时，
\[
Q_5(y)\equiv \prod_{y_i\in\{61,601,740,1696,1946\}}(y-y_i)\pmod{3229},
\]
且勒让德符号满足
\[
\left(\frac{D(61)}{3229}\right)=-1,\qquad
\left(\frac{D(601)}{3229}\right)=+1,\qquad
\left(\frac{D(740)}{3229}\right)=
\left(\frac{D(1696)}{3229}\right)=
\left(\frac{D(1946)}{3229}\right)=-1.
\]
该模式排除“全部 \([\Delta_i]\) 在平方类空间中相等”的退化情形，并与定理 \ref{thm:xi-terminal-zm-delta-node-tangent-wreath} 的 \(\dim_{\FF_2}W=5\) 一致。
\end{remark}

\begin{corollary}[节点分裂统计的显式 Chebotarev 密度]\label{cor:xi-terminal-zm-delta-node-tangent-chebotarev}
令 \(G=\Gal(L/\QQ)\cong C_2\wr S_5\)。
除去有限坏素数后，对素数 \(p\) 定义事件 \(E_p\)：多项式 \(Q_5\) 在 \(\FF_p\) 上完全分裂。
在事件 \(E_p\) 下，五个节点均在 \(\FF_p\) 上有定义；令 \(k(p)\in\{0,1,2,3,4,5\}\) 为其中满足 \(D(y_i)\) 在 \(\FF_p\) 中为平方者的个数（等价地，该节点的两条切线在 \(\FF_p\) 上分裂的个数）。
则对每个 \(k\in\{0,1,2,3,4,5\}\)，满足 \(E_p\) 且 \(k(p)=k\) 的素数集合具有自然密度
\[
\delta_k=\frac{\binom{5}{k}}{|G|}=\frac{\binom{5}{k}}{3840}.
\]
\end{corollary}
\begin{proof}
对 \(L/\QQ\) 应用 Chebotarev 密度定理：任意共轭类 \(C\subset G\) 的素数密度为 \(|C|/|G|\)。
事件 \(E_p\) 等价于 Frobenius 在 \(S_5\) 商上的分量为恒等，因而 Frobenius 落在核 \((C_2)^5\) 中。
在该核内，共轭类由“负号坐标个数”决定；恰有 \(k\) 个坐标取正的 \(S_5\)-轨道大小为 \(\binom{5}{k}\)。
将该共轭类大小代入即得 \(\delta_k=\binom{5}{k}/3840\)。
\end{proof}

\begin{conclusion}[节点分裂方向域与 \texorpdfstring{$\delta$}{delta}-分支域的线性互不相交]\label{con:xi-terminal-zm-delta-node-tangent-b5-linearly-disjoint}
令 \(L/\QQ\) 为定理 \ref{thm:xi-terminal-zm-delta-node-tangent-wreath} 的节点分裂方向域，并令 \(L_B/\QQ\) 为 \(\delta\)-分支五次 \(B_5(\delta)\) 的分裂域（结论 \ref{con:xi-terminal-zm-delta-b5-discriminant-ramification-quadratic}）。
则
\[
L\cap L_B=\QQ,
\qquad
\Gal(LL_B/\QQ)\cong (C_2\wr S_5)\times S_5.
\]
\end{conclusion}
\begin{proof}
由定理 \ref{thm:xi-terminal-zm-delta-node-tangent-wreath}，\(L/\QQ\) 为 Galois 且 \(\Gal(L/\QQ)\cong C_2\wr S_5\)。
由定理 \ref{thm:xi-terminal-zm-delta-a5-b5-galois-s5}，\(L_B/\QQ\) 为 \(S_5\)-伽罗瓦扩张，故其唯一非平凡真正规子域为二次分辨子域
\(\QQ(\sqrt{-140267})\)（结论 \ref{con:xi-terminal-zm-delta-b5-discriminant-ramification-quadratic}）。
若 \(L\cap L_B\neq \QQ\)，则必有 \(\QQ(\sqrt{-140267})\subseteq L\)。
但结论 \ref{con:xi-terminal-zm-delta-node-tangent-quadratic-subfields} 已完全分类 \(L\) 的二次子域，不含 \(\QQ(\sqrt{-140267})\)。
故 \(L\cap L_B=\QQ\)，从而两者线性互不相交，并得到直积伽罗瓦群同构。证毕。
\end{proof}

\begin{conclusion}[分歧源的正交分离：直积伽罗瓦分解下的惯性投影]\label{con:xi-terminal-zm-delta-node-tangent-b5-inertia-orthogonal}
沿用结论 \ref{con:xi-terminal-zm-delta-node-tangent-b5-linearly-disjoint} 的记号，并令
\[
G:=\Gal(LL_B/\QQ)\cong (C_2\wr S_5)\times S_5,
\]
其中两个因子的投影分别记为 \(\mathrm{pr}_1,\mathrm{pr}_2\)。
设 \(p\) 为奇素数且 \(p\neq 3\)。
则：
\begin{enumerate}
  \item 若 \(p\mid \Disc(B_5)\) 且 \(p\nmid \Disc(Q_5)\)，则 \(p\) 在 \(L\) 中不分歧，因而
  \[
  \mathrm{pr}_1\bigl(I_p(LL_B/\QQ)\bigr)=1,\qquad
  \mathrm{pr}_2\bigl(I_p(LL_B/\QQ)\bigr)=I_p(L_B/\QQ).
  \]
  特别地，对 \(p\in\{17,37,223,3739,7333\}\) 成立。
  \item 若 \(p\mid \Disc(Q_5)\) 且 \(p\nmid \Disc(B_5)\)，则 \(p\) 在 \(L_B\) 中不分歧，因而
  \[
  \mathrm{pr}_2\bigl(I_p(LL_B/\QQ)\bigr)=1,\qquad
  \mathrm{pr}_1\bigl(I_p(LL_B/\QQ)\bigr)=I_p(L/\QQ).
  \]
  特别地，对 \(p\in\{11,31,727\}\) 成立。
\end{enumerate}
因此除去 \(\{2,3\}\) 的共同分歧之外，合成扩张 \(LL_B/\QQ\) 在所有奇素数处实现了“节点分裂方向因子”与“分歧值分裂因子”的惯性正交化。
\end{conclusion}
\begin{proof}
由判别式判据，若 \(p\nmid\Disc(f)\) 则由 \(f\) 生成的分裂域在 \(p\) 处不分歧。
对 \(L_B=\mathrm{Spl}(B_5)\) 直接应用结论 \ref{con:xi-terminal-zm-delta-b5-discriminant-ramification-quadratic} 即得其分歧支撑。
\par
对 \(L\)：其底域 \(K=\mathrm{Spl}(Q_5)\) 的分歧支撑包含于 \(\Disc(Q_5)\) 的素因子集合。
而 \(L/K\) 由添入 \(\sqrt{\Delta_i}\) 生成；由命题 \ref{prop:xi-terminal-zm-delta-node-tangent-discriminant-norm-nonsquare} 的范数恒等式 \eqref{eq:xi-terminal-zm-delta-node-tangent-discriminant-norm} 可知每个 \(\Delta_i\) 的有理素因子支撑包含于 \(\{2,3,11,31,727\}\)，故当 \(p\) 为奇素数、\(p\neq 3\) 且 \(p\notin\{11,31,727\}\) 时，\(p\) 在 \(L\) 中不分歧。
\par
最后，由结论 \ref{con:xi-terminal-zm-delta-node-tangent-b5-linearly-disjoint} 的线性互不相交与直积分解，合成域的惯性群在两因子上的投影等于各自子扩张的惯性群；当某一子扩张在 \(p\) 处不分歧时，其惯性为平凡，从而另一因子承载全部惯性。证毕。
\end{proof}

\begin{theorem}[十次层判别式与二次分裂的分歧支撑]\label{thm:xi-terminal-zm-delta-node-preimage-discriminant-and-norm}
令 \(R_{10}(X)\) 如 \eqref{eq:xi-terminal-zm-delta-node-preimage-r10}。
则：
\begin{enumerate}
  \item \(R_{10}\) 在 \(\QQ\) 上不可约。
  \item 其多项式判别式满足严格分解
  \begin{equation}\label{eq:xi-terminal-zm-delta-node-preimage-disc-r10}
  \boxed{\;\Disc(R_{10})=-2^{88}\cdot 3^{17}\cdot 11^{4}\cdot 31^{2}\cdot 727^{2}\cdot 5146068463^{2}\;} .
  \end{equation}
  因而 \(\Disc(R_{10})\) 的平方自由核为 \(-3\)，并且素数 \(37\) 不在判别式支撑中出现。
  \item 令 \(L=\QQ(y)\)（其中 \(y\) 为 \(Q_5\) 的一根），并令 \(K=L(X)\) 由二次中介式 \eqref{eq:xi-terminal-zm-delta-node-preimage-quadratic-mediator} 生成。
  记该二次式对 \(X\) 的判别式为
  \begin{equation}\label{eq:xi-terminal-zm-delta-node-preimage-Delta-y}
  \Delta_X(y)
  :=\bigl(65536y^{4}+47104y^{3}-35392y^{2}-22412y+1774\bigr)^{2}
  -4\cdot 279\cdot\bigl(110592y^{4}+67584y^{3}-60096y^{2}-32496y+3075\bigr).
  \end{equation}
  则其在 \(L/\QQ\) 下的范数满足纯素支撑恒等式
  \begin{equation}\label{eq:xi-terminal-zm-delta-node-preimage-norm-Delta}
  \boxed{\;\mathrm N_{L/\QQ}\!\bigl(\Delta_X(y)\bigr)=-3^{25}\cdot 31^{10}\;} .
  \end{equation}
  从而二次层 \(K/L\) 的分歧仅可能出现在素数 \(3\) 与 \(31\) 上方。
\end{enumerate}
\end{theorem}

\begin{proof}
不可约性可由素数 \(p=17\) 的模约化给出：\(R_{10}\bmod 17\) 在 \(\FF_{17}[X]\) 中保持十次不可约因子，从而由 Gauss 引理推出 \(R_{10}\) 在 \(\QQ\) 上不可约。
\par
判别式分解 \eqref{eq:xi-terminal-zm-delta-node-preimage-disc-r10} 为整系数多项式判别式的直接计算与素因子分解。
平方自由核仅由奇次幂素因子与符号贡献，故为 \(-3\)；判别式支撑中不出现 \(37\) 即说明 \(37\) 不分歧。
\par
对范数恒等式 \eqref{eq:xi-terminal-zm-delta-node-preimage-norm-Delta}，用一般的范数--结式公式（\(\deg\Delta_X=8\) 且 \(Q_5\) 的首项系数为 \(4096\)）：
\[
\mathrm N_{L/\QQ}\!\bigl(\Delta_X(y)\bigr)=4096^{-8}\cdot \Res_{y}\!\bigl(Q_5(y),\Delta_X(y)\bigr).
\]
直接计算结式并约去 \(4096^{8}=2^{96}\) 即得 \eqref{eq:xi-terminal-zm-delta-node-preimage-norm-Delta}。
可复现核验脚本：\texttt{scripts/exp\_xi\_leyang\_node\_preimage\_r10\_wreath\_audit.py}。证毕。
\end{proof}

\begin{theorem}[节点预像十次层的花环积伽罗瓦对称]\label{thm:xi-terminal-zm-delta-node-preimage-galois-wreath}
令 \(G\) 为 \(R_{10}(X)\) 在 \(\QQ\) 上的分裂域的伽罗瓦群，其自然作用于 \(10\) 个根集合。
则
\[
1\longrightarrow (C_2)^5\longrightarrow G\longrightarrow S_5\longrightarrow 1
\]
为自然短正合列，并且该正合列分裂；从而
\[
G\cong C_2\wr S_5\cong (C_2)^{5}\rtimes S_5,
\qquad
|G|=2^{5}\cdot 5!=3840.
\]
等价地，\(G\) 同构于 \(B_5\) 型 Weyl 群（超八面体群），其 \(10\) 点作用为五个二元块上的完全花环作用。
\end{theorem}

\begin{proof}[证明要点]
由定理 \ref{thm:xi-terminal-zm-delta-node-preimage-elimination-r10}，对任一根 \(\alpha\)（\(R_{10}(\alpha)=0\)），存在 \(\beta:=y(\alpha)\) 满足 \(Q_5(\beta)=0\)，且 \(\alpha\) 在 \(\QQ(\beta)\) 上满足二次中介式 \eqref{eq:xi-terminal-zm-delta-node-preimage-quadratic-mediator}。
因此 \([\QQ(\beta):\QQ]=5\)、\([\QQ(\alpha):\QQ(\beta)]=2\)，并且 \(G\) 的 \(10\) 点作用存在大小为 \(2\) 的不变块系（同一 \(\beta\) 之上的两根 \(\alpha\) 形成一块），从而
\[
G\subseteq C_2\wr S_5.
\]
又由定理 \ref{thm:xi-terminal-zm-delta-a5-b5-galois-s5} 可知 \(\Gal(Q_5/\QQ)\cong S_5\)，故块置换商群为 \(S_5\)，并得到短正合列
\[
1\longrightarrow K\longrightarrow G\longrightarrow S_5\longrightarrow 1,
\qquad
K\subseteq (C_2)^5.
\]
\par
由 \eqref{eq:xi-terminal-zm-delta-node-preimage-disc-r10}，\(\Disc(R_{10})\notin\QQ^{\times 2}\)，故 \(G\nsubseteq A_{10}\)。
另一方面，来自块置换的 \(S_5\) 在 \(10\) 点上的提升恒为偶置换，因而奇偶性只能来自核 \(K\) 中的“块内换位”。
这排除 \(K\) 为偶和子群（其仅产生偶置换），从而 \(K\) 仅可能为对角子群或全体 \((C_2)^5\)。
\par
取素数 \(p=101\)。由 \eqref{eq:xi-terminal-zm-delta-node-preimage-disc-r10} 可知 \(101\nmid\Disc(R_{10})\)，故在分裂域中不分歧。
计算 \(R_{10}\bmod 101\) 的不可约因子次数分拆得到 \((1)(1)(8)\)，对应 Frobenius 元在 \(10\) 点上的循环型为 \((1)(1)(8)\)。
而在“仅允许全局块内同时交换”的对角核情形下，任一元素要么为纯块置换（其循环长度在 \(10\) 点上成对出现），要么与全局交换复合（从而不可能保留固定点）。
故循环型 \((1)(1)(8)\) 在该情形下不可能出现，遂排除对角核，必有 \(K=(C_2)^5\)。
综上 \(G=(C_2)^5\rtimes S_5=C_2\wr S_5\)。证毕。
\end{proof}

\begin{corollary}[十次分裂域在节点五次分裂域之上的极大初等 \texorpdfstring{$2$}{2}-扩张结构]\label{cor:xi-terminal-zm-delta-node-preimage-maximal-2-extension}
令 \(L\) 为 \(R_{10}\) 的分裂域，令 \(L_5\) 为 \(Q_5\) 的分裂域。
则 \([L:L_5]=32\) 且 \(\Gal(L/L_5)\cong (C_2)^5\)。
此外，除平凡情形外，\(\Gal(L/L_5)\) 的 \(S_5\)-不变真子模仅有两类：对角线子模与偶奇子模。
因此 \(L/L_5\) 的非平凡正规 \(2\)-子扩张严格只有两条中间正规塔：
\[
L\ \supset\ L^{\mathrm{even}}\ \supset\ L_5,\qquad [L:L^{\mathrm{even}}]=2,\ [L^{\mathrm{even}}:L_5]=16,
\]
以及
\[
L\ \supset\ L^{\mathrm{diag}}\ \supset\ L_5,\qquad [L:L^{\mathrm{diag}}]=16,\ [L^{\mathrm{diag}}:L_5]=2,
\]
其中 \(L^{\mathrm{even}}\) 与 \(L^{\mathrm{diag}}\) 分别对应于“偶数块翻转子群”与“对角线翻转子群”的不动域。
\end{corollary}
\begin{proof}
由定理 \ref{thm:xi-terminal-zm-delta-node-preimage-galois-wreath}，\(\Gal(L/\QQ)\cong (C_2)^5\rtimes S_5\)，并且 \(L_5\) 为块置换商 \(S_5\) 的不动域；因而
\(
\Gal(L/L_5)\cong (C_2)^5
\)
且 \([L:L_5]=2^5=32\)。
将 \((C_2)^5\) 视为 \(\FF_2^5\) 的置换模，则其 \(S_5\)-不变真子模仅为对角线子模与偶奇子模两种非平凡情形。
由 Galois 对应，二者分别对应所示两条中间正规塔。证毕。
\end{proof}

\begin{proposition}[节点预像覆盖的连通性、单轨道性与最小定义域次数]\label{prop:xi-terminal-zm-delta-node-preimage-orbit-degree10}
令 \(S:=\mathrm{Sing}(\overline{\mathcal C}_\delta)\) 为五节点的零维子概形，并令
\[
\Sigma:=\nu^{-1}(S)\subset E(\overline{\QQ})
\]
为归一化预像（按重数计为 \(10\) 点）。
则：
\begin{enumerate}
  \item \(\Sigma\) 在 \(\Gal(\overline{\QQ}/\QQ)\) 作用下构成单一轨道；因而 \(\Sigma(\QQ)=\varnothing\)。
  \item 对任意 \(Q\in \Sigma\)，其最小定义域 \(F=\QQ(Q)\) 满足 \([F:\QQ]=10\)。
  \item \(F\) 的 Galois 闭包之伽罗瓦群同构于 \(C_2\wr S_5\cong W(B_5)\)。
\end{enumerate}
\end{proposition}
\begin{proof}
由定理 \ref{thm:xi-terminal-zm-delta-node-preimage-galois-wreath}，
分裂域的伽罗瓦群同构于 \(C_2\wr S_5\)，其在十次根集合上的标准块作用为传递作用；
故 \(\Sigma\) 为单一伽罗瓦轨道，第一条成立。
\par
对任意点 \(Q\in\Sigma\)，其稳定子同构于 \((C_2)^4\rtimes S_4\)，指数为 \(10\)，从而 \([\QQ(Q):\QQ]=10\)。
并且 Galois 闭包对应于该作用的全体群，故伽罗瓦群同构于 \(C_2\wr S_5\)。
\end{proof}

\endinput


\endinput
